\subsubsection{Reconstruction parallèle du chemin chromosomique}\label{sec:split}
Deux versions d'un assemblage peuvent partager des millions de bases consécutives. Une coupe dans une région identique permet de traiter les deux fragments en parallèle, à condition qu'un chemin optimal passe par cette coupe. La commande \texttt{nvpd-split} propose jusqu'à sept coupes, espacées sur la première chaîne. Elle recherche chaque mot de 16 bases dans une fenêtre de la seconde chaîne, puis vérifie 128 bases autour de la correspondance. Chaque fragment perd son préfixe et son suffixe communs. NvPD traite le rectangle restant jusqu'à 262\,144 cellules, et BiWFA traite les rectangles plus grands~\cite{sadiq,biwfa}.

Si \(C=((i_0,j_0),\ldots,(i_q,j_q))\) contient les extrémités et les coupes proposées, la concaténation des chemins locaux fournit une borne supérieure,
\[
 U=\sum_{k=1}^{q} d(A_{i_{k-1}:i_k},B_{j_{k-1}:j_k})\ge d(A,B).
\]
Les intervalles sont semi-ouverts. Lorsque \(U=|m-n|\), la borne inférieure de longueur suffit. Sinon, tout chemin de coût au plus \(U\) reste dans \(|i-j|\le U\). De chaque côté de la coupe, le programme cherche plus de \(U\) mots disjoints de 16 bases qui n'apparaissent qu'une fois dans la fenêtre accessible de la cible. Une édition rompt au plus un de ces mots disjoints. Un mot subsiste donc de chaque côté sur tout chemin optimal. Entre ces deux correspondances, les chaînes coïncident et la diagonale de coût nul passe par la coupe. Des voisinages disjoints permettent d'appliquer cet argument à toutes les coupes retenues. Un témoin de 16\,384 bases contient deux copies d'un même fragment, distantes de 256 bases. Une substitution efface la première occurrence du mot dans la cible. La seconde propose une coupe dont le décalage imposerait au moins 512 éditions sur les deux fragments. Ce critère écarte cette coupe et le CIGAR global coûte une substitution.

\begin{algorithm}[H]
\caption{Reconstruction par coupes vérifiées}\label{alg:split}
\begin{algorithmic}[1]
\Require Strings \(A,B\), number of pieces \(q\le8\)
\Ensure Optimal global edit score and extended CIGAR
\State \(C\gets\textsc{CandidateCuts}(A,B,q)\)
\State In parallel, \((d_k,P_k)\gets\textsc{AlignTrimmed}(A_k,B_k)\) for each interval of \(C\)
\State \(U\gets\sum_k d_k\), \(V\gets\{c_0,c_q\}\)
\For{each interior \(c_k\in C\)}
  \If{\(U=|m-n|\) or \(\textsc{Witnesses}(A,B,c_k,U)\)} \State \(V\gets V\cup\{c_k\}\) \EndIf
\EndFor
\State Reuse paths of single intervals and realign merged intervals between successive cuts in \(V\)
\State \Return sum of interval scores and concatenated run lengths
\end{algorithmic}
\end{algorithm}

L'algorithme~\ref{alg:split} borne la recherche de témoins à \(U\le1\,024\). Si l'écart de longueur dépasse déjà ce budget, le programme aligne directement le fragment entier après retrait des extrémités communes. Les durées de la figure~\ref{fig:split} couvrent la reconstruction du CIGAR en mémoire. Chaque processus effectue un appel d'échauffement puis un appel mesuré. Le rejeu du CIGAR vérifie ensuite les bases consommées et le coût de toutes les opérations. Pour des coûts de fragments \(t_k\) sur \(p\) travailleurs, le temps parallèle est minoré par \(\max(\max_k t_k,\sum_k t_k/p)\). La vérification des coupes et les intervalles réunis ajoutent du travail séquentiel. Le pic résident inclut le FASTA et le rejeu. Pour le CIGAR d'Edlib sur le chromosome 1, la limite de 120\,s du processus, échauffement compris, a nécessité cinq blocs supplémentaires avec une limite de 720\,s. Chaque bloc relance aussi les trois méthodes concurrentes. Jusqu'à trois blocs de ce complément ont partagé le nœud. Les quatre durées et les quatre RSS du CIGAR du chromosome 1 proviennent de ces mêmes blocs.
\chart{comparison/trace_chromosomes.png}{Le calcul du CIGAR sur le chromosome 1 situe quatre méthodes dans le plan temps--mémoire. Le panneau b) distingue le gain par rapport à la paire entière du gain à découpage fixé. Le dénominateur est le temps des fragments sur huit cœurs. Les symboles donnent les moyennes de cinq blocs, les traits épais un écart-type et les traits fins l’intervalle à 95\,\%.}{fig:split}{bench/cpu/chromosome-trace/measure.py}
Pour le CIGAR du chromosome 1, Edlib prend \(165{,}10\pm0{,}15\)\,s, WFA2 \(183{,}11\pm5{,}14\)\,ms et BiWFA \(381{,}81\pm15{,}28\)\,ms sur un cœur. À huit cœurs, le solveur par fragments prend \(11{,}81\pm1{,}23\)\,ms. Son RSS moyen, FASTA compris, atteint 480,8\,Mio contre 1\,192,5\,Mio pour WFA2.
Six coupes subsistent sur ce chromosome. Le retrait des extrémités communes vide cinq des sept fragments. BiWFA résout les deux autres. Sur les chromosomes 18 et 20, deux petits rectangles par paire appellent NvPD et aucun fragment ne demande BiWFA. Le chromosome 20 conserve six coupes sur sept. À huit cœurs, conserver ces six coupes divise le temps par 28,56 par rapport au rejet de toutes les coupes au premier échec (IC à 95\,\% de 26,76 à 30,36). Son CIGAR prend \(3{,}73\pm0{,}13\)\,ms, contre \(49{,}42\pm0{,}45\)\,ms avec WFA2 et \(39{,}147\pm0{,}030\)\,s avec Edlib sur un cœur.

Dans l'essai où seul le découpage varie, le témoin aligne la paire entière après retrait des mêmes extrémités communes. Sur le chromosome 1, le rapport de son temps à celui de huit travailleurs vaut \(32{,}07\pm0{,}62\). Pour la fenêtre de deux millions de bases entre GRCh38 et CHM13, l'écart de longueur vaut 1\,263 bases et le chemin coûte 5\,599 éditions. Les deux calculs traitent alors un fragment unique. Leur rapport de temps vaut \(1{,}002\pm0{,}017\).

\synthese{Pour deux révisions du chromosome 1, le découpage à huit cœurs reconstruit les 176 éditions en \(11{,}58\pm0{,}14\)\,ms dans l'essai d'ablation. À découpage fixé, le passage d'un à huit cœurs accélère le CIGAR d'un facteur \(1{,}67\pm0{,}04\).}
